\documentclass[11pt,a4paper]{article}

\usepackage[utf8]{inputenc}
\usepackage[indonesian]{babel}
\usepackage{amsmath, amssymb, amsfonts}
\usepackage{geometry}
\geometry{margin=1in, headheight=14pt}
\usepackage{graphicx}
\usepackage{booktabs}
\usepackage{tcolorbox}
\tcbuselibrary{skins,breakable}
\usepackage{enumitem}
\usepackage{fancyhdr}

\definecolor{unibblue}{RGB}{0, 51, 102}
\definecolor{unibgold}{RGB}{212, 175, 55}
\definecolor{darkgreen}{RGB}{34, 139, 34}

\pagestyle{fancy}
\fancyhf{}
\rhead{\textbf{TKE-41XX: Optimisasi untuk Teknik Elektro}}
\lhead{Tugas Rumah / Problem Set Minggu 2}
\rfoot{Halaman \thepage}

\title{\vspace{-1cm}\textbf{PROBLEM SET MINGGU KE-2:}\\Optimisasi Linier (LP) dan Metode Grafis 2D}
\author{Program Studi S1 Teknik Elektro --- Universitas Bengkulu}
\date{Semester Ganjil 2026}

\begin{document}

\maketitle

\begin{tcolorbox}[colback=unibgold!10,colframe=unibgold!80!black,title=\textbf{Petunjuk Pengerjaan dan Tenggat Waktu}]
\begin{itemize}[leftmargin=*]
    \item \textbf{Tenggat Pengumpulan}: 1 Minggu setelah pertemuan kelas.
    \item \textbf{Format Pengumpulan}: Laporan tertulis (PDF) memuat plot grafis manual/komputasi dan skrip SciPy.
    \item **Taksonomi Bloom**: Memuat soal tingkat kognitif **C1 hingga C6**.
\end{itemize}
\end{tcolorbox}

\vspace{0.3cm}

\section*{Soal 1: Konsep Dasar Poligon Feasibel (C1--C2)}

\begin{enumerate}[label=\textbf{1.\arabic*}, leftmargin=*]
    \item \textbf{[C1 --- Mengingat]} Jelaskan definisi formal dari himpunan cembung (*convex set*) dan berikan pembuktian bahwa irisan dari dua himpunan cembung pasti menghasilkan himpunan cembung!
    
    \item \textbf{[C2 --- Memahami]} Jelaskan kondisi matematis yang menyebabkan suatu masalah LP memiliki: (a) Solusi optimal unik, (b) Solusi optimal jamak (*multiple optimal solutions*), dan (c) Tidak terbatas (*unbounded*).
\end{enumerate}

\section*{Soal 2: Solusi Grafis dan Pemrograman Komputasi (C3--C4)}

\begin{enumerate}[label=\textbf{2.\arabic*}, leftmargin=*]
    \item \textbf{[C3 --- Mengaplikasikan]} Diberikan masalah alokasi daya 2 generator linier:
    \[
        \min Z = 20 P_1 + 15 P_2 \quad \text{s.t.} \quad P_1 + P_2 \ge 200, \quad P_1 \le 150, \quad P_2 \le 120, \quad P_1, P_2 \ge 0
    \]
    Gambarkan daerah poligon feasibel secara grafis dan tentukan titik pojok optimal $P_1^*, P_2^*$!

    \item \textbf{[C4 --- Menganalisis]} Tuliskan skrip Python dengan `scipy.optimize.linprog` untuk memverifikasi titik optimal pada Soal 2.1!
\end{enumerate}

\section*{Soal 3: Evaluasi Sensitivitas Grafis dan Desain Model (C5--C6)}

\begin{enumerate}[label=\textbf{3.\arabic*}, leftmargin=*]
    \item \textbf{[C5 --- Mengevaluasi]} Hitunglah rentang koefisien biaya $c_1$ (biaya Generator 1) sedemikian rupa sehingga titik pojok optimal pada Soal 2.1 tetap tidak berubah!

    \item \textbf{[C6 --- Mencipta]} Rancanglah sebuah model LP 2D untuk meminimalkan rugi-rugi transmisi pada 2 saluran kabel sejajar dan selesaikan secara grafis serta komputasional!
\end{enumerate}

\end{document}
